\chapter{Cities Organise Strangers}
\section{An Unremarkable Lump of Clay}
Pick something up. It is so small that your palm barely feels its weight: a clay cone little larger than a peanut, rough, undecorated, dark red after firing. Found in an ancient site's soil, it would probably be discarded as broken brick rubble.

Yet archaeologists have excavated thousands of these cones, balls, and discs in Mesopotamia, the land between the Tigris and Euphrates, roughly today's Iraq. Dating from about eight to five thousand years ago, they are not rubbish. Scholars now broadly identify them as early accounting components, called clay tokens.

Their use ran something like this: owe a storehouse ten sheep and put ten tokens in a bag; borrow a sack of wheat and record another token. Later, hollow clay balls enclosed them. Before sealing, people impressed their shapes on the surface to prevent disputes. Eventually someone recognised that surface marks already recorded the quantity, making enclosed tokens unnecessary: simply inscribe a tablet. Many believe writing grew from accounts this way. Its mother may have been bookkeeping, not poetry.

Hold that clay lump and consider why it exists. Its entire purpose is to settle obligations among people who do not know one another well. Familiar neighbours need none: if Zhang borrows rice from Li, both remember, and so does the village. Only when numbers defeat memory and unfamiliarity defeats trust does fired clay remember for you.

This token is a city's seed. How does the urban machine organise thousands of strangers to specialise, trade, pay taxes, build walls, and dig channels? What does it create, and what does it crush?

\section{A Street Five Thousand Years Ago}
Now enter a scene.

Time: around 5,200 years ago. Place: Uruk in southern Mesopotamia, perhaps already home to tens of thousands. In an age of short lives and mostly hundred-person settlements, it is a megacity.

Sounds fill the street before full daylight. First, donkeys: caravans carrying barley baskets enter the gate, hooves striking packed earth, nostrils steaming. A driver leads them to a large courtyard beside the wall, the temple storage area. A clean-robed scribe sits at its entrance with a half-dry tablet on his knees. Each unloaded basket earns reed-pressed symbols for date, amount, and deliverer's mark. The illiterate driver recognises his own small stone cylinder seal's rolled pattern. He watches the record completed before leaving. The barley is tax, not a gift.

Sun-dried mudbrick houses line the street, mud-plastered layers supporting reed-and-earth roofs. Narrow lanes remain cool in shade until the sun climbs. At one doorway a woman breaks fermented grain cakes into a vat to brew beer. Uruk drinks beer as routinely as you drink milk; workers receive daily beer rations. She brews beyond household needs and trades the excess with neighbours for pots, cloth, or a small bronze knife.

The rising sun brings varied people: reed carriers, boat repairers, wool spinners, temple livestock drivers, northern mountain traders selling stone and timber, speaking different accents. Most cannot name one another; migrants from many villages fill the city. Around noon, temple workers queue for barley rations, their names crossed off tablets. Even the street is an organised project: packed surfaces, roughly directed lanes, and someone investigating walls that encroach on the road.

Above them, an immense central platform carries a white-walled temple visible from afar. It is simultaneously worship site, largest storehouse, largest employer, and largest landholder.

Look around and see another difference from villages: the city's body is artificial and costly throughout. Outer walls of packed earth and mudbrick stand several metres high with watchtowers, consuming thousands of workdays; without them, accumulated wealth waits for raiders. Roads of differing widths carry donkeys and porters and require maintenance. Drains beside lanes and walls carry wastewater downhill; block one and a lane swims in its own filth. Temples and palaces stand tallest. Uruk's temple dominates, while other cities foreground the king's palace, controlling armies, decrees, and tribute. Temple manages divine granaries, palace human swords. Their powers sometimes merge, sometimes contend, while residents in between pay grain to both.

Stay on this street and notice three things. Almost every adult practises a different craft: brewers cannot repair boats, boatwrights cannot keep accounts, scribes cannot farm. A village where everyone farms could scarcely imagine this division of labour. Almost everything moves: grain enters, rations leave, pots, cloth, beer change hands. This is a market. And invisible arrangements sustain movement: accounts, seals, standard measures, and a force making rules count, rather than family affection.

\section{How Can Strangers Live Cooperatively?}
Set out the real problem before answering.

A village of several hundred scarcely needs management. Everyone knows who shirks, defaults on loans, or neglects parents; reputation is a shackle. Hundreds of thousands of years in familiar communities shaped minds to recognise faces and remember favours owed.

Cities shatter that arrangement. On Uruk's street, most daily contacts, bread sellers, grain collectors, ration distributors, lane-sharing neighbours, are strangers briefly met and perhaps never seen again. Reputation, face, kinship, elders' authority: familiar society's mechanisms largely fail.

Then comes the real question:

\textbf{Why can mutually unfamiliar people owing one another nothing share city walls without fighting every day?}

Why should the baker trust tomorrow's repayment of seed grain? Why trust his oil is not watered? If an irrigation channel breaks, why should upstream people repair what benefits downstream fields? Why should families supply labour and food for a hundred thousand cubic metres of wall, protecting households that contribute neither?

This is no question of ancient stupidity. Even today, institutions only partly contain it. A sharper problem crouches behind: when someone offers to organise the solution, how do you know they are not seizing a chance to move everyone's grain into their own house?

Urban history lies between these questions: how strangers cooperate, and whether the power organising them consumes them in return.

\section{Voices Inside and Outside the Walls}
Several positions on this stranger-coordination machine deserve full hearings across history.

\textbf{Voice 1: Storehouses and taxes insure strangers.}

Give temple scribes and storekeepers their full case. Tax collectors are unpopular in every age, but their arguments need not be empty.

Mesopotamian farming depends on two difficult rivers: floods arrive inconveniently, drought can last years. One disaster can ruin a family or village. A storing city saves abundance for scarcity and organises thousands to build canals and embankments beyond any household's means. Standard scales and measures spare strangers repeated quarrels; records and seals make default costly; walls discourage raiders. None is free. Every ration and fair market weighing stands on that collected barley basket. Here tax is an insurance premium, paid normally and claimed during famine or war.

\textbf{Voice 2: Storehouses extract, and walls imprison.}

Give taxed farmers equal seriousness. Their voice is faint in surviving records because writers ate from the stores.

Their account differs. A family works dawn to dusk, then sends much of its harvest into a courtyard it cannot enter. Unreadable tablet marks decide whether enough remains to eat. Failed harvests mean borrowing at strikingly high ancient Near Eastern barley interest; default can cost land, livestock, even children's freedom. Urban laws protect creditors more thoroughly than tenants. Walls cut both ways, excluding enemies and preventing residents' escape. Canal and wall labour demands are vast, without scribes' sons on conscription lists. This city resembles less insurance than a pump extracting rural surplus and ordinary freedom into an organisation.

\textbf{Voice 3: The city is my only chance.}

Between them speaks the voluntary migrant. Many historically chose cities. Village destiny seemed fixed at birth: father's field, neighbouring village's bride, grandfather's fatal illness. Cities offered a new craft and identity. A village shepherd learns seal engraving and owns a shop ten years later; an inheritance-excluded daughter earns wages spinning in a large workshop. Anonymity liberates those seeking changed occupations, changed fortunes, or escape from damaged reputations. Medieval Europe's saying "City air makes you free" referred to serfs attaining freedom after sufficient urban residence. Its survival shows countless people voting with their feet.

All three may describe one city: insurance, extraction, opportunity simultaneously. Which is it for you, and what determines that?

\section[Thinking Bridge: Three Questions]{A Thinking Bridge: Three Questions for Dismantling the Machine}
Any urban arrangement, wall, tax, market, sewer, can be dismantled with three questions. Take this tool beyond antiquity to present institutions, including school.

\textbf{First: What benefits does this arrangement produce, and for whom?}

Importance alone is insufficient; specify benefits and recipients. Walls provide residents safety; standard measures lower traders' transaction costs through fewer arguments and reweighings; storage distributes disaster risk among those able to pay taxes and receive grain. Notice every sentence's second half: benefits never fall evenly.

\textbf{Second: Who bears the costs?}

Publicity most readily conceals this question. Wall accounts list earth and bricks; actual costs include thousands of farmers leaving their work in slack seasons, sometimes busy ones, to ram earth. Standard measures cost almost nothing to maintain, enabling millennia of survival. Stores cost substantial harvest shares paid by cultivators, while administrators take wages first. The grander the praise, the more closely inspect whose bodies support the base.

\textbf{Third: What makes people comply?}

Roughly three things. Voluntary use: everyone likes fair scales that avert disputes. Custom and faith: temple grain is divine offering, refusal disrespect. Coercion: officers and whips enforce refusal. Healthy arrangements often mix all three. Pure coercion is most expensive, requiring many supervisors who themselves need supervising.

Test your familiar institution, school. What benefits whom? Who pays, not merely fees but time and rankings' continual tax on some pupils' self-esteem? Does punctual homework follow choice, custom, or force? An institution's worth depends on the combined shape of these answers, not its slogan.

\section[Two Exhibits: A Stone and a Poem]{Read Two Exhibits: A Stone Pillar and a Poem}
Now primary material.

First, a black stone over two metres tall: Hammurabi's law stele, inscribed in Old Babylon over thirty-seven centuries ago, later looted to Susa in Elam, now in Paris's Louvre. Over two hundred cuneiform provisions concern marriage, loans, employment, medicine, building. Read the building provision, section 229, translated here from the source's Chinese paraphrase:

\begin{quote}
If a builder constructs a house unsoundly, and its collapse kills the owner, the builder shall be put to death.
\end{quote}
The next provision executes the builder's son if the owner's son dies. (Code of Hammurabi, sections 229--230.)

Do not let life-for-life brutality immediately divert you; that deserves another discussion. First inspect the city this rule \textbf{presupposes}. Building is a specialised occupation hired rather than performed by households. Recognised contracts specify parties and responsibility. An authority can execute sentences and expects compliance. Publicly erected stone rules tell strangers consequences in advance. Division of labour, contracts, bureaucracy, force: one short law screenshots the urban operating system.

The second exhibit comes from China. "Unbroken", among the Greater Odes of the Book of Songs, recounts Zhou ancestor Gugong Danfu leading his people below Mount Qi to build a homeland:

\begin{quote}
They summoned the Minister of Works and the Minister of the Multitude to establish dwellings. Their measuring cords ran straight; boards were bound for building; solemn rose the ancestral temple. ("Unbroken", Greater Odes, Book of Songs.)
\end{quote}
Meaning: summon officials responsible for construction and for land and households; build homes, straighten marking lines, bind wall formwork, erect a dignified ancestral temple.

Another screenshot. Before bricks come \textbf{officials}: the Minister of Works directs construction, the Minister of the Multitude people and labour. Organising others is now a profession with specialised posts. The scribes of Mesopotamia and the ministers of Zhou independently reveal the same urban type: people who neither dig nor farm, but record and direct. Scribes, treasurers, works ministers, tax officers: bureaucrats. Cities need them like bodies need nerves; whether those nerves become nutrient-draining tumours is every city's question.

A third, less spectacular but most numerous evidence class is accounts. Most excavated Mesopotamian tablets are neither laws nor epics but lists: barley rations to a work crew on a date, oil from a place received by a store in a year. Too dull to forge for prestige, they become especially reliable testimony. They reveal two things. First, cities operate through accounts rather than households, reducing named people to ration entries and debts. Second, participants understand the system religiously: grain belongs to gods, scribes serve temples, rather than people inhabiting a state or economy. Separate their religious duty from our tax-and-bureaucracy categories to see what happened.

\section{Why Cities Rose: Three Explanations}
Separate four kinds of content. \textbf{What happened}: from over five thousand years ago, Mesopotamia, the Nile, Indus, Yellow and Yangtze river basins, then later the Americas developed large walled, storing, specialised settlements, an evidence-supported fact. \textbf{Participants' understanding}: gods required cities and received grain, their own account. \textbf{Why}: competing explanations. The following three compete to fit facts, not to defend cities.

\textbf{Explanation 1: Surplus nourished cities.}

Most famously systematised as the urban revolution by twentieth-century archaeologist Vere Gordon Childe: agricultural improvement lets one grow food for more than one, freeing artisans, priests, scribes, soldiers to cluster as a city. Neat and intuitive: no spare grain, no street of specialists. But surplus need not become cities. Societies have eaten, shared, or sacrificed it without urbanising. Causation may even reverse: concentrated people and organisation might first force greater surplus.

\textbf{Explanation 2: Great projects forced power into being.}

Canals, embankments, walls require thousands coordinated to one plan and schedule. Controlling coordination controls rations and labour mobilisation; cities crystallise that power. Evidence includes five-thousand-year-old Liangzhu in the Yangtze region, surrounded by multiple dams whose earthworks exceeded any village's labour. Pharaoh-centred Egypt organised stores, surveying, and labour; pyramids monumentally display this capacity. Research describes many builders as ration-paid farmers serving in slack seasons rather than legendary slaves, further highlighting organisation. Its weakness: it explains wall bricks better than people's hearts. If cities only coerced, why did so many willingly risk migrating into them?

\textbf{Explanation 3: Faith gathered people first.}

Temples may precede cities: repeated ritual and festival gatherings become settlement. American examples lend force. Mexico's Teotihuacan, flourishing around fifteen centuries ago with over a hundred thousand residents, follows a broad avenue bordered at its ends and side by immense pyramids. The plan itself is religious. Mississippi Valley Cahokia, around a millennium ago, held over ten thousand residents and great flat-topped ritual mounds. Attraction seems initially rooted in access to gods. Yet faith cannot explain daily food and wastewater after gathering: it explains the heart, not the gut.

Each grasps truth: surplus supplies capital, projects power, faith willing arrivals. Unequal but not wholly conflicting, they require caution against any claim of completeness. Explaining an institution's origin is not defending it. Irrigation generating central authority describes causation, not justice or inevitability. An Indus city next challenges that supposed inevitability.

\section[Further Challenge: Sewer Politics]{A Deeper Challenge: The Politics of Sewers}
Add two economic concepts, intimidating names with simple meanings.

First, \textbf{public goods}. Two features define them: excluding non-payers is difficult, as walls protect untaxed residents too; another beneficiary does not reduce others' share, as my streetlight also illuminates you. Clean streets, drainage, walls, and security approximate public goods. Trouble follows: if nonpayment still benefits me, why not let others build? This is \textbf{free-riding}. Universal free-riding leaves no ride: unrepaired canals, unbuilt walls, sewage in streets. Almost every core urban facility must clear this hurdle.

See its reality at Mohenjo-daro in today's Pakistan, around forty-five centuries ago, perhaps tens of thousands strong at its height. Its most astonishing feature is not a palace but \textbf{drainage}. Almost every house has bathing facilities and a well. Wall-side pottery pipes and brick channels feed covered street drains, interconnected at precise gradients with periodic settling pits for cleaning. This was the world's most sophisticated urban sanitation, unmatched in many cities for millennia afterward.

Here lies our \textbf{counterexample}, deliberately breaking prejudice. You may have absorbed the formula city equals temple plus palace plus king, great works equal strongman's compulsion. Yet excavation reveals no generally accepted palace, lavish royal tombs, or kingly statues trumpeting victories. Instead: ordered blocks, a great public bath, strikingly standard bricks, and drains. Numerous inscribed Indus seals remain undeciphered, leaving even rulers' names unknown.

Its force lies in showing that stranger cooperation need not have one answer. No visible king, yet maintained citywide sewers. Merchant associations, religious communities, some unimaginable negotiation mechanism? Honestly, we do not know. But apparently obvious claims that large cooperation requires autocracy acquire a crack.

How did civilisations handle free-riding? Several mixed prescriptions recur. Coercion: taxes and labour backed by whipping and imprisonment, prominent in Mesopotamia and Egypt. Faith: divine works and meritorious service, common to temples and festivals. Expanded reciprocity: small communities extend turn-taking labour across canal-building alliances. Markets: charge where possible, as at ferries and toll bridges. None cures everything, each costs something: coercion enlarges power, faith grants priests interpretive control, reciprocity limits scale, markets exclude those unable to pay.

A slightly melancholy conclusion: public goods work best when \textbf{invisible}. Nobody remembers functioning drains; only blockages, floods, and disasters make news. Most governance achievements appear as nothing happening. This makes judging cooperation machines difficult: benefits stay silent, costs speak loudly.

\section[Sewers Below, Addresses in Your Phone]{The Sewers Below You, the Addresses in Your Phone}
These questions did not remain ancient.

Look down. Flushing, drawing tap water, walking on unsinking pavement: each descends directly from Mohenjo-daro's channels. You barely notice payment, though household utilities and taxes provide it. Invisibility marks good functioning. Three days without water would instantly reveal that ancient engineering's greatness.

Painful tuition bought the system. Nineteenth-century London was the world's largest city yet suffered repeated deadly cholera. In 1854 John Snow mapped Soho deaths and found clustering around the Broad Street public pump. He persuaded authorities to remove its handle, and the outbreak subsided. This became a classic epidemiological and public-health moment: a city first located disease through maps and numbers rather than prayer and expulsion. Notice its structure. Snow's map and Uruk's tablet perform the same work: collecting scattered strangers' facts so the machine knows where repair is needed. The ledger now records causes of death rather than debts.

Look up. China's present urbanisation, history's largest, moves hundreds of millions from villages within decades, magnifying Uruk's street a hundred thousand times. Old recipes return. A delivered meal coordinates cook, rider, platform, and many strangers through accounts. Standards become ratings and algorithms, five stars or reduced visibility as digital fair scales. Accounts become phone payments: clay tokens grow into digits you still trust.

Old problems remain too. Cities produce opportunities and hierarchies: strong schools, jobs, hospitals alongside expensive housing and fierce competition. Migrants face ancient thresholds in household registration, school districts, eligibility, invisible walls replacing brick. Density multiplies prosperity and risk: close contact spreads ideas, wealth, disease. Cities amplified epidemics; fourteenth-century Black Death followed trading cities and killed roughly a third to half of Europe. A recent global pandemic again showed that a machine gathering millions also efficiently passes viruses among them. Pollution enters the same ledger: factories and cars discharge costs into shared air, charged even to non-drivers.

Young people's dinner tables still debate big-city migration. Go: opportunity, experience, freedom, liberating city air. Stay: prices, commutes, loneliness, the extracting pump. You now know this argument is at least five thousand years old, with evidence on both sides.

\section{Your Turn: Two Sides of a Clay Token}
Return to the clay lump.

One side says accounting: strangers can calculate, trust, cooperate. Without it, no stores, walls, drains, or familiar modern comforts.

The other says power: recordkeepers know everyone's circumstances; storekeepers hold everyone's food. Accounts are neutral tools and early nerves of rule. The line from tablets to platform data never broke; reeds merely became servers.

Your question follows:

\textbf{If a machine offers safety, order, convenience, and opportunity in exchange for recording and ranking you and taking some earnings, while you can never be certain it takes only a reasonable share, would you live inside or move somewhere without it?}

Give both sides their full reasons. For entry: outside, disasters and bandits are free too, and nobody builds your sewers. For leaving: benefits stay silent, but the machine's hand speaks; finer ledgers bring recordkeepers closer. Then inspect your reasoning: do you fear coercion or insufficient provision? Love opportunity or the machine's very existence?

Uruk's donkey-driving taxpayer had no choice five millennia ago. Today you do. The token reaches your palm. How will you spend it?
